
\documentclass[12pt,a4paper]{article}

\usepackage[margin=1in]{geometry}
\usepackage{graphicx}
\usepackage{xcolor}
\usepackage{array}
\usepackage{float}
\usepackage{caption}

\definecolor{custompurple}{RGB}{200,20,200}
\colorlet{customorange}{red!60!yellow}

\title{Assignment 04 - Lists, Colors, Figures etc.}
\author{Manasa Paulson}
\date{\today}

\begin{document}
	
	\maketitle
	\tableofcontents
	\newpage
	
	\section{Lists}
	
	Following are some technologies that can be used to format content:
	
	\begin{itemize}
		\item \LaTeX
		\item HTML (+CSS)
		\item Markdown
	\end{itemize}
	
	You would most likely use \LaTeX{} in the following scenarios:
	
	\begin{enumerate}
		\item To write research papers.
		\item To embed math equations within other software
		(e.g. GDocs, HTML etc.).
		\item To create your CV, project reports and presentations.
	\end{enumerate}
	
	There are three main categories of electronic markup:
	
	\begin{description}
		\item[Presentational] Used by traditional word-processing systems:
		binary codes embedded within document text that produce the
		WYSIWYG (what you see is what you get) effect.
		e.g. Microsoft Word.
		
		\item[Procedural] Markup is embedded in text which provides
		instructions for programs to process the text. e.g. TeX.
		
		\item[Descriptive] Markup is specifically used to label parts
		of the document for what they are, rather than how they should
		be processed. e.g. \LaTeX{}, HTML5.
	\end{description}
	
	\section{Tables and Colors}
	
	Code for Table 1 is written in this subsection. However, it will
	be rendered on a special page. The table is positioned at the
	top of a page using the optional float argument.
	
	\begin{table}[!t]
		\centering
		\caption{Comparison of different technologies}
		\label{tab:technologies}
		
		\begin{tabular}{|p{2.5cm}|p{2cm}|p{2cm}|p{3.5cm}|}
			\hline
			Language & Requires compilation & Output & Use \\
			\hline
			C/C++ & Yes & Executable & To make apps \\
			\hline
			\LaTeX & Yes & PDF & To make documents \\
			\hline
			HTML/CSS/JS & No & Webpage & To make websites \\
			\hline
			\multicolumn{1}{|c|}{2.5 cm} &
			\multicolumn{1}{c|}{2 cm} &
			\multicolumn{1}{c|}{2 cm} &
			\multicolumn{1}{c|}{3.5 cm} \\
			\hline
		\end{tabular}
	\end{table}
	
	Colors can be represented using RGB values or hexadecimal
	notation in HTML.
	
	A custom color is defined using 200 red, 20 green and 200 blue:
	
	\begin{verbatim}
		\definecolor{custompurple}{RGB}{200,20,200}
	\end{verbatim}
	
	Another custom color is defined using 60\% red and 40\% yellow:
	
	\begin{verbatim}
		\colorlet{customorange}{red!60!yellow}
	\end{verbatim}
	
	The following equation uses the custom colors:
	
	\[
	\textcolor{custompurple}{T(n)}
	=
	\textcolor{customorange}{aT(n/b)}
	+
	\textcolor{custompurple}{f(n)}
	\]
	
	\section{Figures}
	
	\subsection{Pictures: Rotation and Cropping}
	
	The MNIT logo is rotated using the graphicx package.
	
	\begin{figure}[H]
		\centering
		\rotatebox{90}{
			\includegraphics[width=0.25\textwidth]
			{"C:/Users/user/Downloads/mnitlogo.jpg"}
		}
		\caption{MNIT logo}
		\label{fig:logo}
	\end{figure}
	
	The MNIT building image is cropped using the trim and clip
	options. The top portion of the image is cropped to remove
	the sky, and the image is fitted to the text width.
	
	\begin{figure}[H]
		\centering
		\includegraphics[
		width=\textwidth,
		trim=0 0 0 100,
		clip
		]{"C:/Users/user/Downloads/building.jpg"}
		\caption{This is a floating image and can appear anywhere.}
		\label{fig:building}
	\end{figure}
	
	\subsection{Multiple Images}
	
	This subsection includes two minipages with equal width.
	
	\begin{figure}[H]
		\centering
		
		\begin{minipage}{0.45\textwidth}
			\centering
			\includegraphics[width=\textwidth]
			{"C:/Users/user/Downloads/mnitlogo.jpg"}
			\captionof{figure}{This is the first minipage}
		\end{minipage}
		\hfill
		\begin{minipage}{0.45\textwidth}
			\centering
			\includegraphics[width=\textwidth]
			{"C:/Users/user/Downloads/building.jpg"}
			\captionof{figure}{This is the second minipage}
		\end{minipage}
		
	\end{figure}
	
\end{document}
